% Generated draft: requires mathematical and source review.
\documentclass[11pt]{article}
\usepackage[margin=25mm]{geometry}
\usepackage{fontspec}
\setmainfont{Latin Modern Roman}
\usepackage{amsmath,amssymb,mathtools,unicode-math}
\setmathfont{Latin Modern Math}
\usepackage{enumitem,needspace,xurl,hyperref}
\setlength{\parindent}{0pt}
\setlength{\parskip}{5pt}
\setlength{\emergencystretch}{3em}
\newcommand{\N}{\mathbb{N}}
\newcommand{\Z}{\mathbb{Z}}
\newcommand{\Q}{\mathbb{Q}}
\newcommand{\R}{\mathbb{R}}
\newcommand{\C}{\mathbb{C}}
\newcommand{\F}{\mathbb{F}}
\newcommand{\A}{\mathbb{A}}
\newcommand{\PP}{\mathbb{P}}
\newcommand{\E}{\mathbb{E}}
\newcommand{\eps}{\varepsilon}
\newcommand{\dd}{\,\mathrm{d}}
\begin{document}
\section*{UnsolvedMath ID: 800840}
\textbf{Source problem code:} catalogue entry
\subsection{Definitions and mathematical statement}
Let \(d\in\N\) tend to infinity, and let \(n=n(d)\in\N\) be a sequence of positive integers. Write \(I_d\) for the \(d\times d\) identity matrix. For each \(d\), the source specifies random vectors \(x_1,\ldots,x_{n(d)}\in\R^d\) by
\[
x_i\sim\mathcal N(0,I_d/d),\qquad 1\leq i\leq n(d).
\]
The displayed sampling specification gives Gaussian marginal distributions but does not explicitly state independence. The intended joint sampling law must therefore be clarified.

Write \(A_d\) for the event described in the source as the existence of ``an ellipsoid fit to \((x_i)_{i=1}^{n(d)}\).'' The supplied record does not define this phrase. In particular, it does not specify whether the ellipsoid is centered at the origin, whether fitting requires every point to lie on its boundary, or whether degenerate quadratic forms are permitted. Consequently, \(A_d\) remains an unresolved source predicate, and the following is a threshold statement conditional on clarifying that predicate and the joint sampling law.

\textbf{Conjecture 5 (supplied formulation).} For every \(\varepsilon>0\) and every positive integer sequence \(n(d)\),
\[
\limsup_{d\to\infty}\frac{n(d)}{d^2}\leq\frac{1-\varepsilon}{4}
\quad\Longrightarrow\quad
\lim_{d\to\infty}\mathbb P(A_d)=1,
\]
and
\[
\liminf_{d\to\infty}\frac{n(d)}{d^2}\geq\frac{1+\varepsilon}{4}
\quad\Longrightarrow\quad
\lim_{d\to\infty}\mathbb P(A_d)=0.
\]
The source supplies no conclusion at the critical ratio \(1/4\). Its registry status is \texttt{open}; this is supplied classification, not an independently verified mathematical or literature assessment.
\subsection{Short English statement}
The supplied conjecture predicts an ellipsoid-fitting threshold at \(n/d^2=1/4\) for Gaussian points: fitting should occur with probability tending to one below the stated lower threshold and with probability tending to zero above the stated upper threshold. The meaning of fitting and the intended independence assumption require clarification before this becomes a fully specified mathematical statement.
\subsection{Sources}
\begin{itemize}
\item UnsolvedMath local registry, permanent ID \(800840\), exact source problem code \texttt{catalogue entry}. The supplied record labels the problem \texttt{open} and attributes it to Bandeira et al., with proposed year 2025.
\item Bandeira et al., \emph{Randomstrasse101: Open Problems of 2024} (2025), Conjecture 5, context: ellipsoid fitting. Supplied source URL: \url{https://arxiv.org/abs/2504.20539}. The record reports extraction by \texttt{arxiv\_latex} and an access date of 2026-10-16. These are supplied provenance details; the source was not independently visited or verified for this draft.
\end{itemize}
\end{document}
